% © 2026 Ing. David Jaroš — CC BY-NC-ND 4.0
%
% This work is licensed under a Creative Commons Attribution-NonCommercial-NoDerivatives
% 4.0 International License (CC BY-NC-ND 4.0).
%
% License History: Earlier drafts (up to v0.3) were released under CC BY 4.0.
% From v0.4 onward, all material is released under CC BY-NC-ND 4.0 to protect
% the integrity of the theoretical work during ongoing academic development.
%
% See LICENSE.md for full license text.

\ifdefined\INCLUDEMODE
\else
  \documentclass[12pt,a4paper]{article}
  \usepackage[a4paper, margin=2.5cm]{geometry}
  \usepackage{amsmath, amssymb, amsthm}
  \usepackage{mathtools}
  \usepackage{hyperref}
  \usepackage{xcolor}
  \usepackage{mdframed}
  \usepackage{booktabs}
  \usepackage{array}

  \newtheorem{theorem}{Theorem}[section]
  \newtheorem{lemma}[theorem]{Lemma}
  \newtheorem{proposition}[theorem]{Proposition}
  \newtheorem{corollary}[theorem]{Corollary}
  \theoremstyle{definition}
  \newtheorem{definition}[theorem]{Definition}
  \theoremstyle{remark}
  \newtheorem{remark}[theorem]{Remark}

  \newmdenv[backgroundcolor=yellow!20, linecolor=orange, linewidth=1.5pt]{gapbox}
  \newmdenv[backgroundcolor=green!15, linecolor=green!60!black, linewidth=1.5pt]{resultbox}
  \newmdenv[backgroundcolor=blue!8, linecolor=blue!50, linewidth=1.5pt]{insightbox}

  \title{\textbf{Discrete Symmetries in Biquaternion Field Theory:\\
                 Formal Involutions, $C$, $P$, $T$, and $CPT$}}
  \author{Ing.\ David Jaro\v{s}}
  \date{2026}

  \begin{document}
  \maketitle

  \begin{abstract}
  We give a self-contained, reviewer-grade treatment of the discrete symmetries
  $C$ (charge conjugation), $P$ (spatial parity), $T$ (time reversal), and
  $CPT$ within the biquaternion algebra $\mathcal{B}=\mathbb{C}\otimes\mathbb{H}$
  and for the fundamental UBT field $\Theta(q,\tau)$ over complex time
  $\tau = t + i\psi$.
  In Section~\ref{sec:ds:algebra} we catalogue all canonical involutions
  of $\mathcal{B}$ — complex conjugation, quaternion conjugation, Hermitian
  conjugation, reversion, and the grade involution — and identify their physical
  roles.
  In Section~\ref{sec:ds:CPT} we give the canonical definitions of $C$, $P$,
  $T_{\mathrm{UBT}}$, and their composition $CPT$, expressed as algebra
  involutions acting simultaneously on field values and spacetime arguments.
  Section~\ref{sec:ds:tables} collects operator and notation tables for
  reference.
  Section~\ref{sec:ds:noncomm} discusses noncommutativity constraints.
  Conventions follow \texttt{canonical/algebra/involutions\_Z2xZ2xZ2.tex};
  earlier detailed derivations are in
  \texttt{symmetry/step1\_CPT\_definitions.tex}.
  \end{abstract}
\fi

%──────────────────────────────────────────────────────────────────────────────
\section{The Biquaternion Algebra and Its Involutions}
\label{sec:ds:algebra}
%──────────────────────────────────────────────────────────────────────────────

\subsection{Algebra Setup}

The biquaternion algebra is
\begin{equation}
  \mathcal{B} = \mathbb{C} \otimes_{\mathbb{R}} \mathbb{H},
  \label{eq:ds:B}
\end{equation}
an associative, non-commutative $\mathbb{R}$-algebra of dimension $8$.
A basis element is written $c\,\mathbf{e}$ with $c \in \mathbb{C}$ and
$\mathbf{e} \in \{1,\mathbf{I},\mathbf{J},\mathbf{K}\}$.
A general element reads
\begin{equation}
  b = b_0 + b_1\mathbf{I} + b_2\mathbf{J} + b_3\mathbf{K},
  \quad b_a \in \mathbb{C}.
  \label{eq:ds:general}
\end{equation}
The \emph{central} imaginary unit $i\in\mathbb{C}$ commutes with all
$\mathbf{I},\mathbf{J},\mathbf{K}$.  The quaternion units satisfy
$\mathbf{I}^2=\mathbf{J}^2=\mathbf{K}^2=-1$, $\mathbf{I}\mathbf{J}=\mathbf{K}$,
$\mathbf{J}\mathbf{K}=\mathbf{I}$, $\mathbf{K}\mathbf{I}=\mathbf{J}$, and are
non-commutative: $\mathbf{I}\mathbf{J}=-\mathbf{J}\mathbf{I}$.

\subsection{Five Canonical Involutions}
\label{sec:ds:involutions}

An \emph{involution} on $\mathcal{B}$ is an $\mathbb{R}$-linear map
$\sigma:\mathcal{B}\to\mathcal{B}$ with $\sigma^2=\mathrm{id}$.
We list the five physically relevant involutions.

\subsubsection{$P_1$: Complex Conjugation}
\begin{equation}
  P_1(b) = \bar{b} := \bar{b}_0 + \bar{b}_1\mathbf{I}
            + \bar{b}_2\mathbf{J} + \bar{b}_3\mathbf{K},
  \label{eq:ds:P1}
\end{equation}
i.e., $i\mapsto -i$ while $\mathbf{I},\mathbf{J},\mathbf{K}$ are fixed.
$P_1$ is a \emph{ring anti-automorphism restricted to the $\mathbb{C}$ factor};
since $i$ is central, it extends to an $\mathbb{R}$-algebra automorphism
of $\mathcal{B}$.
\textbf{Physical role}: charge conjugation (flips $\mathbb{C}$ phase; see
Section~\ref{sec:ds:CPT}).

\subsubsection{$P_2$: Quaternion Conjugation}
\begin{equation}
  P_2(b) = \tilde{b} := b_0 - b_1\mathbf{I} - b_2\mathbf{J} - b_3\mathbf{K},
  \label{eq:ds:P2}
\end{equation}
i.e., $\mathbf{I},\mathbf{J},\mathbf{K}\mapsto -\mathbf{I},-\mathbf{J},-\mathbf{K}$
while $i$ (and hence $b_a$) is fixed.
$P_2$ is the standard quaternion conjugation, extended $\mathbb{C}$-linearly.
$P_2$ is a ring anti-automorphism: $P_2(bc)=P_2(c)P_2(b)$.
\textbf{Physical role}: spatial parity (flips vector part; see
Section~\ref{sec:ds:CPT}).

\subsubsection{$P_{12}$: Hermitian Conjugation (Full Conjugation)}
\begin{equation}
  P_{12}(b) = b^\dagger := P_1(P_2(b)) = P_2(P_1(b))
  = \bar{b}_0 - \bar{b}_1\mathbf{I} - \bar{b}_2\mathbf{J} - \bar{b}_3\mathbf{K}.
  \label{eq:ds:P12}
\end{equation}
Since $P_1$ and $P_2$ commute ($P_1 P_2 = P_2 P_1$ on $\mathcal{B}$),
$P_{12} = P_1 \circ P_2$ is itself an involution.
It is the standard Hermitian adjoint on $\mathcal{B}$ when the latter is
identified with $M_{2}(\mathbb{C})$ via the standard representation.
\textbf{Physical role}: $CPT$ at the algebra level; adjoint for inner products.

\subsubsection{$\rho$: Reversion}

Reversion is defined on a product of generators by reversing their order:
\begin{equation}
  \rho(\mathbf{e}_{a_1}\cdots\mathbf{e}_{a_k})
  = \mathbf{e}_{a_k}\cdots\mathbf{e}_{a_1}.
  \label{eq:ds:reversion}
\end{equation}
On single generators: $\rho(\mathbf{I})=\mathbf{I}$, $\rho(i)=i$.
On degree-2 products: $\rho(\mathbf{I}\mathbf{J}) = \mathbf{J}\mathbf{I} = -\mathbf{K}$.
For a general element \eqref{eq:ds:general}, reversion fixes $b_0$ and
the vector part $(b_1,b_2,b_3)$ (since each basis vector is grade-1 and
$\rho(\mathbf{e})=\mathbf{e}$):
\begin{equation}
  \rho(b) = b_0 + b_1\mathbf{I} + b_2\mathbf{J} + b_3\mathbf{K} = b.
  \label{eq:ds:reversion_on_B}
\end{equation}
Thus reversion acts as the identity on $\mathcal{B}$
(the $\mathbb{H}$ part is generated by degree-1 elements and does not
generate independent degree-2 components in our basis \eqref{eq:ds:general}).
Reversion becomes non-trivial on tensor products involving Clifford algebra
$\mathrm{Cl}(n)$ embedded in $\mathcal{B}$; in that setting it distinguishes
bivector from vector parts.
\textbf{Physical role}: In extensions of UBT to full Clifford algebras,
reversion provides a natural CPT-like operator; in the current
$\mathcal{B}$-only framework it coincides with the identity.

\subsubsection{$\pi$: Grade (Parity) Involution}

The grade involution is the $\mathbb{Z}_2$ grading map:
\begin{equation}
  \pi(b) = b_0 - b_1\mathbf{I} - b_2\mathbf{J} - b_3\mathbf{K} = P_2(b).
  \label{eq:ds:grade}
\end{equation}
Grade-even elements (scalars $b_0$) are $\pi$-eigenvectors with eigenvalue $+1$;
grade-odd elements (vectors $b_1\mathbf{I}+\cdots$) with eigenvalue $-1$.
In the present $\mathcal{B}$ structure, the grade involution coincides
with $P_2$.
\textbf{Physical role}: separates scalar (spin-0) from vector (spin-1) sectors
of $\Theta$; eigenspaces provide the \emph{UBT chirality analogues} for
left/right projections (candidate mapping to Weyl chirality; proof via
spinor/matrix representation required;
see \texttt{chirality\_and\_parity\_breaking.tex}).

\subsection{Involution Summary}
\label{sec:ds:involution_summary}

\begin{table}[h]
\centering
\caption{Canonical involutions of $\mathcal{B}=\mathbb{C}\otimes\mathbb{H}$.
  Notation: $b = b_0 + b_1\mathbf{I} + b_2\mathbf{J} + b_3\mathbf{K}$,
  $b_a\in\mathbb{C}$.}
\label{tab:ds:involutions}
\begin{tabular}{@{}lllll@{}}
\toprule
Involution & Symbol & Action on $b$ & Type & Physical role \\
\midrule
Complex conjugation   & $P_1$ & $b_a \mapsto \bar{b}_a$ (all $a$)
  & $\mathbb{R}$-alg.\ auto & $C$ \\
Quaternion conjugation & $P_2$ & $\mathbf{e}_i \mapsto -\mathbf{e}_i$ ($i=1,2,3$)
  & ring anti-auto & $P$, grade \\
Hermitian conjugation & $P_{12}$ & $b_a \mapsto \bar{b}_a$, $\mathbf{e}_i\mapsto -\mathbf{e}_i$
  & ring anti-auto & $CP$, adjoint \\
Reversion             & $\rho$   & $b \mapsto b$ (identity in $\mathcal{B}$)
  & ring auto & trivial in $\mathcal{B}$ \\
Grade involution      & $\pi$    & $= P_2$ in $\mathcal{B}$
  & ring auto & chiral split \\
Axis-flip ($\mathbf{I}$) & $P_3$ & $\mathbf{J},\mathbf{K}\mapsto -\mathbf{J},-\mathbf{K}$
  & inner auto & isospin sub-symm. \\
\bottomrule
\end{tabular}
\end{table}

\begin{remark}[Commutativity of $P_1, P_2$]
$P_1$ and $P_2$ commute: $P_1 P_2 = P_2 P_1$.  Proof: both act on different
factors of the tensor product ($\mathbb{C}$ and $\mathbb{H}$ respectively).
Thus the group generated by $\{P_1,P_2\}$ is $\mathbb{Z}_2\times\mathbb{Z}_2$,
with the four elements $\{\mathrm{id}, P_1, P_2, P_1P_2\}$ corresponding
to $\{1, C, P, CP\}$ at the algebra level.
\end{remark}

\begin{remark}[Noncommutativity and $C$, $P$]
$C$ and $P$ commute \emph{as operators on field values} because $P_1 P_2 = P_2 P_1$.
They do not generally commute when their \emph{coordinate actions} are included,
but the coordinate actions of $C$ ($\tau\to\tau^*$) and $P$ ($q\to\bar{q}$)
also commute.  Hence $[C,P]=0$ as operators on the field space.
\end{remark}

%──────────────────────────────────────────────────────────────────────────────
\section{Canonical Definitions of $C$, $P$, $T$, $CPT$}
\label{sec:ds:CPT}
%──────────────────────────────────────────────────────────────────────────────

The fundamental field is $\Theta(q,\tau)\in\mathcal{B}$, with
$q = q^0 + \mathbf{q} = q^0 + q^1\mathbf{I}+q^2\mathbf{J}+q^3\mathbf{K}\in\mathcal{B}$
encoding spacetime position and $\tau = t + i\psi\in\mathbb{C}$ the complex
time ($t\in\mathbb{R}$, $\psi\in\mathbb{R}$ the auxiliary imaginary-time
coordinate).

\begin{insightbox}
\textbf{Convention: $q^0$ vs.\ $\tau$ (physical time vs.\ internal scalar coordinate).}
\begin{itemize}
  \item $\tau = t + i\psi$ is the \textbf{physical complex time parameter}.
        Its real part $t\in\mathbb{R}$ is the observable (physical) time.
        Its imaginary part $\psi\in\mathbb{R}$ is the auxiliary imaginary-time
        coordinate on the $\psi$-circle.
  \item $q^0\in\mathbb{R}$ is the \textbf{scalar (temporal) part of the
        biquaternion coordinate} $q\in\mathcal{B}$.  In the standard
        identification $q \leftrightarrow x^\mu$, one sets
        $q^0 = x^0 = ct$, so that $q^0$ coincides with the real-time
        component.  However, $q^0$ is an \emph{internal} scalar coordinate
        of the algebra element $q$, not an independent time variable.
  \item \textbf{Double-counting warning}: if $\tau$ already encodes physical
        time $t$, one must not simultaneously treat $q^0$ as an independent
        physical time.  The canonical convention is: $\tau$ is the physical
        time parameter; $q^0 = t$ is its embedding in $q$ as the scalar
        component.  They are the same degree of freedom, not two separate times.
  \item \textbf{Transformation rules}: under $T_{\mathrm{UBT}}$ ($\tau\to-\tau$),
        both $t\to-t$ and $\psi\to-\psi$.  Consequently $q^0\to-q^0$ as well
        (since $q^0=t$).  This must be implemented consistently in every CPT
        analysis.
\end{itemize}
\end{insightbox}

\begin{definition}[Charge Conjugation $C$]
\label{def:ds:C}
\begin{align}
  C[\Theta](q,\tau) &:= P_1\!\bigl(\Theta(q,\tau)\bigr), \label{eq:ds:C_field} \\
  C: \tau &\mapsto \tau^* = t - i\psi. \label{eq:ds:C_tau}
\end{align}
$C$ complex-conjugates all biquaternion components while leaving the spacetime
coordinate $q$ unchanged.
For gauge fields: $C: A_\mu^a T_a \mapsto -A_\mu^a T_a^*$
(charge reversal; $U(1)$ charges flip sign).
\end{definition}

\begin{definition}[Spatial Parity $P$]
\label{def:ds:P}
\begin{align}
  P[\Theta](q,\tau) &:= P_2\!\bigl(\Theta(\bar{q},\tau)\bigr), \label{eq:ds:P_field} \\
  P: q &\mapsto \bar{q} = q^0 - \mathbf{q}, \label{eq:ds:P_coord} \\
  P: \tau &\mapsto \tau \quad \text{(parity-even)}. \label{eq:ds:P_tau}
\end{align}
Spatial inversion $(t,\mathbf{x})\mapsto(t,-\mathbf{x})$ is encoded in
$q\mapsto\bar{q}$; the algebra value is simultaneously transformed by $P_2$.
For gauge fields: $P: A_0(x)\mapsto A_0(\bar{x})$, $P: A_i(x)\mapsto -A_i(\bar{x})$.
\end{definition}

\begin{definition}[Time Reversal $T_{\mathrm{UBT}}$]
\label{def:ds:T}
The canonical UBT time reversal is \emph{full complex time negation}:
\begin{align}
  T_{\mathrm{UBT}}[\Theta](q,\tau) &:= \Theta(q,-\tau), \label{eq:ds:T_field} \\
  T_{\mathrm{UBT}}: \tau &\mapsto -\tau = -t - i\psi. \label{eq:ds:T_tau}
\end{align}
No algebra involution is applied to the field value.
For gauge fields: $T: A_0(t,\mathbf{x})\mapsto A_0(-t,\mathbf{x})$,
$T: A_i(t,\mathbf{x})\mapsto -A_i(-t,\mathbf{x})$.
\end{definition}

\begin{remark}[Antiunitary $T$ vs.\ $T_{\mathrm{UBT}}$]
\label{rem:ds:antiunitary}
In standard QFT, time reversal is antiunitary: $T_{\rm std}[\Theta]=P_1(\Theta(q,-\tau^*))$,
i.e., $t\to-t$ combined with complex conjugation of the field.
The UBT choice $T_{\mathrm{UBT}}$ is \emph{unitary} (no $P_1$) but sends
$\psi\to-\psi$ as well as $t\to-t$.
Whether $T_{\mathrm{UBT}}$ should instead be composed with $P_1$ (yielding
antiunitary $T_3 = P_1 \circ T_{\mathrm{UBT}}$) depends on the canonical
inner product chosen for $\Theta$ modes.
\textbf{Open Problem}: fix the canonical sesquilinear form and determine
the correct antiunitary structure.
See \texttt{open\_problems.md}, item OP-S1.
\end{remark}

\begin{definition}[$CPT$ Operator]
\label{def:ds:CPT}
\begin{equation}
  CPT := C \circ P \circ T_{\mathrm{UBT}}.
  \label{eq:ds:CPT}
\end{equation}
Explicitly:
\begin{equation}
  CPT[\Theta](q,\tau)
  = (P_1\circ P_2)\bigl(\Theta(\bar{q},-\tau)\bigr).
  \label{eq:ds:CPT_explicit}
\end{equation}
In the real-time limit $\psi\to0$:
$(q,\tau)\to(-q^0-\mathbf{q},-t) = $ full four-vector inversion,
reproducing the standard QFT $CPT$ transformation.
\end{definition}

\begin{proposition}[$CPT$ reduces to standard QFT in the $\psi\to0$ limit]
\label{prop:ds:CPT_limit}
Setting $\psi=0$ (hence $\tau=t$, $\tau^*=t$, $-\tau=-t$):
\begin{equation}
  CPT[\Theta]\bigl((q^0,\mathbf{q}),t\bigr)
  = (P_1\circ P_2)\bigl(\Theta(q^0 - \mathbf{q}, -t)\bigr),
\end{equation}
which is the standard $CPT$ operation: inversion of all four spacetime
coordinates combined with Hermitian conjugation of the internal algebra value.
\end{proposition}
\begin{proof}
At $\psi=0$: $\tau=-\tau$ becomes $t\to-t$; $q\to\bar{q}$ flips spatial
components; $(P_1\circ P_2)$ conjugates the field value.
This matches the standard CPT: $(x^\mu)\to(-x^\mu)$ with conjugation.
\end{proof}

%──────────────────────────────────────────────────────────────────────────────
\section{Operator and Notation Tables}
\label{sec:ds:tables}
%──────────────────────────────────────────────────────────────────────────────

\begin{table}[h]
\centering
\caption{Operator table: discrete symmetry operators in UBT.}
\label{tab:ds:operators}
\small
\begin{tabular}{@{}p{2cm}p{4cm}p{4.5cm}p{3.5cm}@{}}
\toprule
Operator & Algebraic action on $\Theta$ value & Field argument action & Physical interpretation \\
\midrule
$C$ & $P_1$: $b_a \mapsto \bar{b}_a$ & $\tau \mapsto \tau^*$ & Charge conjugation; matter $\leftrightarrow$ antimatter \\
$P$ & $P_2$: $\mathbf{e}_i \mapsto -\mathbf{e}_i$ & $q \mapsto \bar{q}$ & Spatial parity; mirror reflection \\
$T_{\mathrm{UBT}}$ & identity & $\tau \mapsto -\tau$ & Full complex time reversal \\
$CP$ & $P_{12}= P_1 P_2$: both & $q\mapsto\bar{q}$, $\tau\mapsto\tau^*$ & Combined charge+parity \\
$CPT$ & $P_{12}= P_1 P_2$ & $q\mapsto\bar{q}$, $\tau\mapsto-\tau$ & Full discrete symmetry; reduces to std.\ $CPT$ at $\psi=0$ \\
$P_\psi$ & identity & $\psi\mapsto-\psi$ & Imaginary-time parity; $n \leftrightarrow -n$ winding modes \\
$P_3$ & $\mathbf{J},\mathbf{K}\mapsto-\mathbf{J},-\mathbf{K}$ & none & Axis-flip; isospin sub-symmetry \\
\bottomrule
\end{tabular}
\end{table}

\begin{table}[h]
\centering
\caption{Transformation of basic fields under $C$, $P$, $T$.
  Sign $+$ = invariant; sign $-$ = odd; $*$ = complex conjugated.}
\label{tab:ds:fieldtransf}
\begin{tabular}{@{}lcccc@{}}
\toprule
Field & $C$ & $P$ & $T_{\mathrm{UBT}}$ & $CPT$ \\
\midrule
$\Theta$ (scalar part $b_0$)     & $b_0^*$ & $+$ & $+$ & $b_0^*$ \\
$\Theta$ (vector part $b_i$)     & $b_i^*$ & $-$ & $+$ & $-b_i^*$ \\
$A_0$ (gauge, temporal)          & $-A_0$  & $+$ & $+$ & $-A_0$ \\
$A_i$ (gauge, spatial)           & $-A_i$  & $-$ & $-$ & $+A_i$ (sign from $P$ and $T$) \\
$F_{0i}$ (electric component)    & $-F_{0i}$ & $-$ & $+$ & $+F_{0i}$ \\
$F_{ij}$ (magnetic component)    & $-F_{ij}$ & $+$ & $+$ & $-F_{ij}$ \\
$F_{\mu\nu}F^{\mu\nu}$           & $+$     & $+$ & $+$ & $+$ \\
$F_{\mu\nu}\tilde{F}^{\mu\nu}$ (dual) & $+$ & $-$ & $-$ & $+$ \\
$\tau$                           & $\tau^*$ & $+$ & $-\tau$ & $-\tau^*\to-\tau$ at $\psi=0$ \\
\bottomrule
\end{tabular}
\end{table}

%──────────────────────────────────────────────────────────────────────────────
\section{Noncommutativity Constraints}
\label{sec:ds:noncomm}
%──────────────────────────────────────────────────────────────────────────────

The non-commutativity of $\mathcal{B}$ has two relevant consequences for
discrete symmetries:

\begin{enumerate}
  \item \textbf{Order matters in products.}
    When $\Theta$ appears in a Lagrangian term as $\Theta^a \Phi^b$ for
    non-central $a,b$, applying $P_2$ gives
    $P_2(\Theta^a\Phi^b) = P_2(\Phi^b)P_2(\Theta^a)$
    (anti-automorphism), not $P_2(\Theta^a)P_2(\Phi^b)$.
    Action terms must be constructed as real traces or norms to remain
    real-valued and well-defined under $P$:
    $\mathrm{Tr}[\Phi^\dagger\Psi] = \mathrm{Tr}[P_2(\Psi)P_2(\Phi)^\dagger]$.

  \item \textbf{Noncommutativity does not affect $C$.}
    $P_1$ acts on the complex scalar factor and commutes with quaternion
    multiplication; thus $P_1(\Theta\Phi)=P_1(\Theta)P_1(\Phi)$
    (algebra automorphism).  Charge conjugation is \emph{not} sensitive
    to ordering.

  \item \textbf{Left-vs-right action.}
    When $\Theta$ transforms under a gauge group $G$ as
    $\Theta \to U\Theta V^\dagger$ (bimodule), the $P$ and $C$ actions
    must specify which side receives the algebra involution.
    The canonical convention is: both $P_2$ and $P_1$ act \emph{simultaneously}
    on left and right, since they are algebra automorphisms/anti-automorphisms
    of $\mathcal{B}$ itself.
\end{enumerate}

\begin{gapbox}
\textbf{Open: Canonical inner product and antiunitary $T$.}
The antiunitary structure of $T$ requires specifying the canonical
sesquilinear form $\langle\cdot,\cdot\rangle$ on the space of $\Theta$ fields.
Until this is fixed, the canonical status of $T_{\mathrm{UBT}}$ as the unique
time reversal remains a model choice rather than a theorem.
See \texttt{open\_problems.md}, item OP-S1.
\end{gapbox}

\ifdefined\INCLUDEMODE
\else
  \end{document}
\fi
